% Generated by `taskgen course examples` from lattices.babai_hard (seed 4).
% Do not edit: change examples.yaml and regenerate.
\begin{example}[алгоритм Бабаи, плохой базис]\label{ex:lattice-babai-bad}
В решётке $\mathcal{L}$, заданной базисными векторами \[v_1 = \begin{pmatrix}-41 \\ 35\end{pmatrix}, v_2 = \begin{pmatrix}87 \\ -82\end{pmatrix},\] найти вектор, ближайший к вектору $w = (-39, 15)^T$.

\par\noindent\textit{Решение.} $\det \mathcal{L} = 317$, $HR(\mathcal{L}) = \left(\frac{|\det\mathcal{L}|}{\|v_1\|\|v_2\|}\right)^{1/2} = 0.222$ --- базис плохой. Алгоритм Бабаи на плохом базисе даёт далёкий вектор (для исходного базиса получилось бы $(-72, 46)^T$, $\|w - r\| \approx 45.277$), поэтому сначала редуцируем базис. \par Алгоритм Гаусса: если $\|v_2\| < \|v_1\|$ "--- меняем местами (${}^{\leftrightarrow}$); $m = \lfloor v_1 \cdot v_2 / \|v_1\|^2 \rceil$; если $m = 0$ "--- стоп, иначе $v_2 \gets v_2 - m v_1$.
\[\begin{array}{c|r|r|r|r|r}
 & v_1 & v_2 & \|v_1\|^2 & v_1\cdot v_2 & m \\ \hline
1 & (-41, 35) & (87, -82) & 2906 & -6437 & -2 \\
2^{\leftrightarrow} & (5, -12) & (-41, 35) & 169 & -625 & -4 \\
3 & (5, -12) & (-21, -13) & 169 & 51 & 0 \\
\end{array}\]
Редуцированный базис $u_1 = (5, -12)^T$, $u_2 = (-21, -13)^T$, $HR = 0.994$. Решаем $w = t_1 u_1 + t_2 u_2$ (правило Крамера): $t_1 = \frac{-822}{317} \approx -2.593$, $t_2 = \frac{393}{317} \approx 1.240$; округляем: $-3 u_1 + 1 u_2 = (-36, 23)^T$, $\|w - r\| \approx 8.544$.\[r = (-36, 23)^T,\quad HR(\mathcal{L}) = 0.222.\]
\end{example}
